\chapter{综合讨论}\label{chap:discussion}

\section{三类证据链（静态--动态--形式化）的互证关系}
静态分析回答“改了什么结构”，动态分析回答“运行时发生了什么”，形式化分析回答“是否仍存在可达反例”。三者构成互补证据链，提升结论可信度。

\section{高风险模块治理启示（storage/replication/WAL）}
针对高风险模块（\texttt{storage}/\texttt{replication}/\texttt{WAL}），建议建立分层治理机制：高频提交监测、关键路径断言增强、回归用例优先配置以及形式化检查前置。

\section{方法局限与有效性威胁}
本文方法存在若干局限：样本时间窗口有限、标注仍受人工判断影响、形式化模型存在抽象损失。外部有效性方面，不同项目的开发流程差异可能影响结论迁移。